% !TeX root = ../papers/nc-p3.tex
\section{種数三の二重被覆のノルムと八基点の母線への集中}
\label{sec:sectional-hyperelliptic-norm}
\begin{theorem}\label{thm:sectional-hyperelliptic-norm}
十八次の双有理二次核曲面モデルで、逆系が $S_0=\mathbb F_0$ または $\mathbb F_2$ 上の
$A=2L+2F$ で、一般切断 $J$ が自己準同型曲線に同型とする。ここで $L$ は核超平面、
$F$ は母線である。このとき長さ八の基点スキーム $Z$ は一般 $J$ と
四母線の和 $\pi^*D_4$ の完全交叉である。$D_4$ は重複を許す固定された次数四の因子である。
\end{theorem}
\begin{proof}[一般切断の判別式]
入力切断の四空間を $V$ とし、入力点の空間は $\mathbb P(V^*)$ と書く。
$Z$ は局所完全交叉で長さ $A^2-8=8$ である。実際一般 $J$ 上で入力射が正則なので
そのイデアルは $J$ 上で可逆、$S_0$ 上で二生成である。
$J$ は母線底への有限平坦二重被覆で $\omega_J=2F|_J$、入力束は
$H|_J=(A|_J)(-Z)$。定理\ref{thm:sectional-cofactor-length}から $2L-H=\omega_J$ が
$J$ 上で成立し、従って $\mathcal O_J(Z)=\mathcal O_J(4F)$。
底の非零な四次式で $Z$ 上消えるものがあれば、その一般 $J$ 上の零因子は
長さ八ゆえ $Z$ に一致し、結論となる。以下そのような四次式がないと仮定する。

$a\in V$ に対応する一般切断 $J$ の相対微分 $w_a=d_{\pi}a|_J$ は
$(A+\omega_\pi)|_J=4F|_J$ の切断である。局所的な母線二次式の微分であり、
その平方は判別式 $\Delta_a$ の引戻し、二重被覆の共役で符号を変える。
有限二重被覆の跡分解（算術種数三より非自明な直和因子は $\mathcal O(-4)$）で、
$w_a$ は $H^0(J,4F)$ の底から来ない一次元部分を張る。
よって $Z$ を切る切断はスカラーを選んで $p_a+w_a$ と書ける。
$p_a$ は底の四次式で、$Z$ への制限の単射性から一意であり $a$ に線形に依存する。
実際 $Z$ 上では全ての $a$ が消えるため $d_\pi a|_Z$ は共通の束の切断として線形である。
ノルムの因子は固定サイクル $\pi_*[Z]$ であり、これを切る八次式を $h$ とすると
\[
           p_a(t)^2-\Delta_a(t)=h(t) C(a)
\]
と書ける。$C$ は $V$ 上の非零な二次式である。一般 $a$ のノルムは非零で、
任意の $a$ への式の延長は係数の多項式性による。

一般の $t$ を固定し、$a$ の母線上の三係数を $A_a,B_a,C_a$ とする。
これらは独立である。一般母線の入力像は次数二の双有理な滑らかな円錐曲線なので完全二次系である。
\[
  a\longmapsto
  \begin{pmatrix} A_a &(B_a+p_a)/2\\(B_a-p_a)/2&C_a\end{pmatrix}
\]
の行列式は $h(t)C(a)/4$ である。$C$ が階数四ならこの線形写像は同型であり、
母線上の二次評価は二行行列に対する階数一の線形汎関数であるから、全ての入力円錐曲線は
固定双対二次曲面 $C^\vee$ に入る。整八次に反する。
階数三ならこの線形写像は階数三で、その核は $C$ の固定根基そのもの。
全 $t$ の母線二次式が零になる非零な入力切断ができてしまう。
$\Delta$ の階数は三なので $C$ の階数は二未満にはならない。
残る階数二では $C=c d$ が独立な二一次式の積である。
この二入力点 $c,d$ は各一般入力円錐曲線上にある。
実際写像の核を除いた三次元空間で判別式 $B^2-4AC$ を使えば、
$p^2-(B^2-4AC)$ が階数二になるのは $p$ の双対判別形式でのノルムが
$B$ のそれに等しい場合である（階数一更新の行列式公式）。
二変数の基底を替えて $p=\pm B$ とでき、二つの因子は $A,C$、
すなわち二つの相異なる点での二次評価である。
\end{proof}

\begin{proof}[二つの固定入力点を通る円錐曲線族]
定理\ref{thm:cofactor-eighteen-birational-quadric}および
\ref{thm:cofactor-eighteen-cone-not-birational}の相対完全系を使う。
すなわち $W=\pi_*(I_Z A)$ は非負直和次数、階数三、次数四で、
$\Gamma=\operatorname{Bl}_Z S_0$ は $\mathbb P(W)$ に写り、像 $Y$ は $2\xi$ の整因子である。
ここで $\xi$ は相対超平面束で、入力束 $H$ に引き戻される。
$c,d$ 上の水平曲線は入力で縮約され、完全な $|H|$ でも縮約される。
$\xi$ の正次数の直和座標が一つでも非零なら完全系から底座標が復元できる。
従って相異なる二水平縮約は二つの自明直和因子を要し $W=\mathcal O^2\oplus\mathcal O(4)$。
相対平面の座標を $u,v,w$ とすれば、$w=0$ 上で $Y$ の式は二つの異なる定点を
切る定数二次式である（恒等零なら $Y$ がこの部分スクロールになる）。
従って四次式倍の $w$ で $u,v$ をずらすと
$Y=(uv=b_8(t)w^2)$、$b_8\ne0$ と書ける。これは特異点が高々有限の整超曲面で正規である。
$\Gamma\to Y$ は双有理で、有限となる開集合では同型となる。

入力四座標のうち二つを $w r_1(t),w r_2(t)$、残る二つを $u+w b_1(t),v+w b_2(t)$ と書ける。
$r_i,b_i$ は四次式で $r_i$ は独立である。固定二点の直線を $\Lambda$ とし、
$r_1,r_2$ の共通因子を除いた底の射の次数を $k\le4$ とする。
一般入力母線平面は $\Lambda$ を通り、平面束の方向はこの底の射で与えられる。
共通零点の上の平面は入力で $\Lambda$ に写る。$\Lambda$ の外では残る各母線平面が
その入力平面に同型に写る。従って $\Gamma$ から入力への射が有限である $\Lambda$ 外では、
繊維は底の次数 $k$ の因子の部分スキームである。ここでは $M_D=f_*L$ で、
これは同じ二次曲面モデルの $\chi(\Gamma,L+nH)=4(n+1)^2$ と最大Cohen--Macaulay性による。
階数一の印は従って長さ三を要する。

垂直な正次元入力繊維の像は高々三点である。実際縮約される母線は
$\ell(Z\cap F)=2$ のものだけで、四本あれば $Z$ は四母線の和にスキームとして含まれる。
これは底の四次式の制限が単射との仮定に反する。基点上の例外直線は入力に縮約されない。
また水平縮約の入力は $c,d$ のみである。従って $\Lambda$ の外の有限な印は
少なくとも $20-3-3=14$。これで $k\le2$ は不可能である。
\end{proof}

\begin{proof}[次数三・四の平面束]
平面の入射三次元多様体を $\Sigma\subset\mathbb P^1\times\mathbb P^3$ とする。
その式 $g_x(t)$ は二次数 $(k,1)$ で共通因子を除いており、$\Sigma$ は射影平面束である。
$\Gamma$ の底と入力の対による像 $Y'$ は因子で、
その類は $(8-2k,2)$。実際一般平面繊維の像は二次曲線、入力超平面の三乗は
$\Sigma$ 上で $k$、$Y'$ における自乗は八である。
$\Lambda$ 外では先の平面座標の同型から $Y'$ と $Y$ は同一視できる
（$r_1,r_2$ の共通零点の入力像は $\Lambda$ 内で、射は固有）。従ってそこで
入力有限なら長さ三の条件は $Y'$ のスキーム繊維そのものである。

$k=3$ なら $Y'$ の式は周囲の $(2,2)$ 式に持ち上がる。
$\Sigma$ の式で割る完全列の左端は $\mathcal O(-1,1)$ である。
長さ三条件には、入力点でその二項二次式が三次の繊維全体で消えることが必要であり、
三つの入力二次係数は全て零になる。係数に共通因子はない。
あればその入力曲面の逆像が $Y'$ に含まれ、整な $Y'$ の入力像が八次であることに反する。
その包の一般二式は互いに素で、定理\ref{thm:sectional-nineteen-stein-cubics}の
三次評価による共通印上界は十二。十四印に届かない。

$k=4$ では $Y'=Y$ であり、$\Lambda$ 上の入力像は $c,d$ だけなので有限である印は
少なくとも $20-2-3=15$ 個で、全て $\Lambda$ 外に取れる。
$Y$ の式 $q\in H^0(\Sigma,\mathcal O(0,2))$ の境界は
$H^1(\mathbb P^1,\mathcal O(-4))\otimes H^0(\mathbb P^3,\mathcal O(1))$
に入る。これは二項二次式上の一次汎関数 $l(x)$ で入力に線形である。
各 $x\notin\Lambda$ に制限すると、四点因子 $D_x=(g_x=0)$ 上での $q(x)$ の、
定数を核とする境界写像になる。
長さ三繊維では $q(x)\mathcal O_{D_x}$ は長さ一、従って $l(x)$ は零ではなく
その台の一点における二次評価に比例する。実際境界のSerre双対の対は
$D_x$ の随伴跡に $q(x)$ と二項二次式の積を送るものなので、
この点で消える全二次式を消す。定数の非零元は長さ一のイデアルを生成しない。
従って十五印が二次Veronese曲線の入力線形引戻しに載る。
非零二次式の上界は十四印である。引戻しが恒等零なら線形像は高々固定一点の二次評価であり、
全ての印で同じ $t$ が $D_x$ に入る。これは一つの入力平面に十五印を強制する。
以上で底の四次式が存在しない場合を全て排除した。
\end{proof}
\path{computations/verify_sectional_hyperelliptic_norm.py} はノルムの行列式、
二点を与える分解および三次・四次平面束の印数を検算する。
